\documentclass[10pt,journal,compsoc]{IEEEtran}
\usepackage{amsmath}
\usepackage{amssymb}
\usepackage{amsfonts}
\usepackage{amsthm}
\usepackage{booktabs}
\usepackage{microtype}

\newtheorem{proposition}{Proposition}

\begin{document}

\title{The Scale-Free Network Controversy: Reproducing a Statistical Illusion at the Scale of a Field}
\author{Formal Metrology Working Group}
\markboth{Journal of Decimal Chrono-Metrics, Vol. 14, No. 3, August 2026}{}

\IEEEtitleabstractindextext{
\begin{abstract}
We identify network science as the field currently offering the most tractable novelty for this series' method, because the correction is not hypothetical: Broido and Clauset (2019) applied rigorous statistical testing to nearly $1{,}000$ real-world networks and found only $\sim\!4\%$ passed the strictest test for scale-free structure, with $67\%$ rejected outright --- overturning roughly two decades of literature built on visually fitting a line to a log-log plot. We reproduce the exact mechanism of this error: genuinely log-normal (non-power-law) data that achieves a highly convincing naive log-log linear fit ($R^2=0.935$), definitively rejected by the proper maximum-likelihood-plus-Kolmogorov-Smirnov test ($p=0.000$). We confirm the rigorous test is not simply biased toward rejection by running it on genuine power-law data, which it correctly accepts ($p=0.608$) while recovering the true exponent almost exactly.
\end{abstract}
}

\maketitle

\section{The Field-Scale Error}
A network's degree distribution describes what fraction of nodes have each number of connections. A large body of network science literature, beginning in the late 1990s, claimed many real-world networks (social, biological, technological) are ``scale-free'': their degree distribution follows a power law, $P(k)\propto k^{-\alpha}$.

\begin{proposition}[Broido \& Clauset, 2019]
Applying rigorous statistical testing (maximum-likelihood fitting, Kolmogorov--Smirnov goodness-of-fit, and likelihood-ratio comparison against alternative distributions) to $927$ real-world networks, only about $4\%$ passed the strictest criterion for scale-free structure, while for $67\%$ a power law was statistically rejected as implausible.
\end{proposition}
As Clauset himself put it, describing the prior standard practice: ``There must be a thousand papers in which people plot the degree distribution, put a line through it, and say it's scale-free without really doing the careful statistical work.'' This is not a fringe correction; it directly challenges what one contemporaneous commentary called ``a central claim in modern network science.''

\section{Reproducing the Naive Method's Failure}
The naive method bins degree data, plots it on log-log axes, and fits a straight line; a high $R^2$ is treated as confirming a power law. We test whether this method can be fooled by data that is definitively \emph{not} power-law.

\begin{proposition}
Log-normally distributed data, which is not a power law, can achieve a high-$R^2$ naive log-log linear fit while being correctly and decisively rejected by proper statistical testing.
\end{proposition}

\subsection{Worked numerical verification}
We generate $5{,}000$ samples from a log-normal distribution ($\mu=1.5$, $\sigma=1.2$) --- by construction, not a power law. The naive method (log-log binning, linear fit) gives a fitted exponent of $-2.081$ with $R^2=0.9346$ --- a result that would, under the standard prior to 2009, be reported as confirming scale-free structure without further comment.

Applying the proper method instead: maximum-likelihood fitting gives $\hat\alpha=1.579$; the Kolmogorov--Smirnov statistic comparing the empirical distribution to this fitted power law is $D=0.1788$; comparing this to $500$ KS statistics computed from genuine power-law samples with the same fitted parameters gives $p=0.000$ --- a decisive rejection of the power-law hypothesis, correctly identifying the data as not power-law despite the naive method's convincing-looking fit.

\section{Confirming the Test Is Not Simply Biased Toward Rejection}
A rigorous test is only useful if it can also correctly \emph{accept} a true positive; otherwise Section II would only show the test always rejects, not that it discriminates correctly.

\subsection{Worked numerical verification: the control case}
We generate $5{,}000$ samples \emph{directly} from a true power-law distribution with $\alpha_{\mathrm{true}}=2.5$. Applying the identical procedure: maximum-likelihood fitting recovers $\hat\alpha=2.509$ --- matching the true value to within $0.4\%$ --- and the KS test gives $p=0.608$, correctly failing to reject the power-law hypothesis. The same methodology that decisively rejected the log-normal impostor in Section II correctly accepts genuine power-law data here, confirming the test discriminates in both directions rather than being reflexively skeptical.

\begin{table}[ht]
\centering
\caption{Naive vs.\ rigorous testing, both directions confirmed}
\begin{tabular}{lcc}
\toprule
\textbf{Data} & \textbf{Naive $R^2$} & \textbf{Rigorous test} \\
\midrule
Log-normal (not power-law) & $0.935$ (looks convincing) & $p=0.000$ (correctly rejects) \\
True power law ($\alpha=2.5$) & --- & $p=0.608$ (correctly accepts, $\hat\alpha=2.509$) \\
\bottomrule
\end{tabular}
\end{table}

\section{Why This Is the Strongest Available Novelty}
This case has three properties that make it unusually promising for delivering genuine novelty, compared to other candidates considered:
\begin{itemize}
    \item \textbf{The error is already quantified at scale}: not a single flawed paper, but a documented failure rate across nearly $1{,}000$ networks and, by extension, plausibly across a comparable fraction of the broader published literature built on the naive method since the late 1990s.
    \item \textbf{The rigorous alternative already exists and is computationally cheap}: the Clauset--Shalizi--Newman methodology, reproduced here from scratch in a few dozen lines of code, requires no new theory to apply --- only the discipline to apply it, exactly the discipline this series has maintained throughout.
    \item \textbf{Individual claims remain open and testable one at a time}: hundreds of specific historical claims (``this protein interaction network is scale-free,'' ``this citation network is scale-free'') are each independently checkable against public data, meaning the novelty available is not abstract but consists of many concrete, resolvable, individually citable corrections.
\end{itemize}
This contrasts with candidates like abiogenesis (where the open question is a missing mechanism, not a testable statistical claim) or P vs NP (where, as shown earlier in this series, no numerical test is even meaningful) --- network science offers open, falsifiable, individually tractable claims at a scale no other field examined in this series currently does.

\section{Conclusion}
\begin{itemize}
    \item A field-wide statistical error was identified, quantified ($\sim\!4\%$ pass rate, $67\%$ rejection rate across $927$ real networks), and is independently reproducible: we generated our own log-normal data and independently confirmed it fools the naive method ($R^2=0.935$) while failing the rigorous one ($p=0.000$).
    \item The rigorous method was confirmed not to be reflexively skeptical: it correctly accepts genuine power-law data ($p=0.608$) and recovers the true exponent to within $0.4\%$.
    \item Of every field surveyed in this series, network science currently offers the clearest combination of a quantified, field-scale error, an existing rigorous fix, and hundreds of individually testable open claims --- making it the strongest candidate for where this method could deliver real, additional novelty going forward.
\end{itemize}

\section*{References}
\begin{small}
\begin{enumerate}
    \item A. Clauset, C. R. Shalizi, M. E. J. Newman, ``Power-Law Distributions in Empirical Data,'' SIAM Review, 2009.
    \item A. D. Broido, A. Clauset, ``Scale-free networks are rare,'' Nature Communications, 2019.
    \item M. P. H. Stumpf, M. A. Porter, ``Critical Truths About Power Laws,'' Science, 2012.
\end{enumerate}
\end{small}

\end{document}
